Questo elaborato definisce la progettazione e la descrizione di un protocollo di sicurezza per la realizzazione di un referendum binario digitale. Il protocollo da realizzare sarà utilizzato nell'ambito di votazioni universitarie, di conseguenza non dovrà essere eseguito in contesti particolarmente critici come delle votazioni nazionali. Quindi, alcune vulnerabilità o limitazioni non ammissibili in contesti particolarmente complessi saranno considerate accettabili nel contesto universitario scelto. \\ \\
L'elaborato è suddiviso in quattro sezioni. Il \textbf{Work Package 1} descrive in generale il modello del sistema, definendo gli attori coinvolti, i possibili avversari e le proprietà che si vogliono rispettare. Il \textbf{Work Package 2} mostra il protocollo progettato partendo dal modello definito al Work Package 1. Nel \textbf{Work Package 3} si analizza la sicurezza del protocollo proposto e infine il \textbf{Work Package 4} spiega i dettagli implementativi del protocollo progettato.